\section{Lightcone construction}
\label{sec:lightcone}

Full-sky catalogues are built from a single simulation volume observed from a
box corner: with the observer at a corner of the (periodic) box, one octant of
the sky is covered out to the comoving distance of the far corner, and eight
runs with the observer at the eight corners of the \emph{same realization}
tile the full sky. For the Websky-configuration box ($L=5236\Mpch = 7.7$~Gpc,
Section~\ref{sec:valcat}) an octant reaches $z_{\rm max}=4.5$ without
replication. This is the construction of the original Websky release; the
radial-shell tiling used by more recent efforts to reach higher redshift is
not needed at this depth.

In lightcone mode, the collapse condition is evaluated at the lookback epoch
of each candidate: a proto-halo at comoving distance $\chi$ from the observer
must collapse by $z(\chi)$ rather than $z=0$, which raises the effective
linear threshold by the inverse growth factor. \ppjl\ evaluates the full
ellipsoidal-collapse table at each candidate's epoch. Displacements and
velocities (equation~\ref{eq:2lpt}) are likewise evaluated at $z(\chi)$, with
the Eulerian position solved consistently with the lightcone crossing.
\needcheck{one-paragraph description of the exact chi-iteration used in
merge/finalize --- verify against current code path before submission (plan
item 3.7)}

Each halo record carries the Lagrangian centre, Eulerian position, velocity,
$R_{\rm TH}$, and (optionally) the strain eigenvalues and second-order terms
used by downstream painting (``extended'' output). Abundance matching
(Section~\ref{sec:algorithm-2lpt}) is applied per octant with the observer's
per-axis position defining $z(\chi)$ for the redshift binning.
